% !TEX spellcheck = en_US
% !TEX spellcheck = LaTeX
\documentclass[12pt,aspectratio=169]{beamer}
\input{header_beamer}
\usetheme{Boadilla}
\usecolortheme{beaver}
\setbeamertemplate{page number in head/foot}{}
\setbeamertemplate{navigation symbols}{}

\title{Lecture-01: Sample and Event Space}
\date{}

\begin{document}
\pagestyle{empty}
\maketitle

\begin{frame}{Notation}

    \begin{itemize}
        \item<1-> Set of first $N$ positive integers: $[N] \triangleq \set{1, 2, \dots, N}$
        \item<2-> Set of positive integers (natural numbers): $\N \triangleq \set{1, 2, \dots}$
        \item<3-> Set of integers: $\Z \triangleq \set{\dots, -1, 0, 1, \dots}$
        \item<4-> Set of non-negative integers: $\Z_+ \triangleq \set{0, 1, \dots}$
        \item<5-> Set of rational numbers: $\Q$, % \triangleq \set{\frac{p}{q}: p, q \in \Z\text{ and are coprime }, q \neq 0}$
        \item<6-> Set of reals: $\R$
        \item<7-> Set of non-negative reals: $\R_+$
        \item<8-> Power set of a set $A$: $\cP(A) \triangleq \set{B: B \subseteq A}$
    \end{itemize}
\end{frame}

\begin{frame}{Function}
    \begin{block}{Definition}<1->
        For sets $A, B$, we denote a function from set $B$ to set $A$ by $f: B \to A$, where $(b, f(b)) \in B \times A$ and for each $b \in B$ there is only one value $f(b) \in A$. That is, \EQ{\set{(b, f(b)): b \in B} \subseteq B \times A.}
    \end{block}
    \begin{itemize}
        \item<2-> Set $B$ is called the \textbf{domain} of function $f$
        \item<3-> Set $A$ is called \textbf{co-domain} of function $f$
        \item<4-> Set $f(B) =\set{f(b) \in A: b \in B}$ is called the \textbf{range} of function $f$
    \end{itemize}
\end{frame}

\begin{frame}
    \begin{example}
        \begin{itemize}
            \item Let $B = \set{1,2,3}$ and $A = \set{a,b}$ 
            \item <2->$\set{(1,a), (2,a), (3,b)}$ corresponds to a function $f:B \to A$ such that $f(1)= f(2) = a, f(3) =b$.
            \item<3-> We will denote this function by ordered tuple $(aab)$.
        \end{itemize}
    \end{example}
 
    \begin{block}{Notation}<4-> 
        The collection of all $A$-valued functions with the domain $B$ is denoted by $A^B$. 
    \end{block}
    
    \begin{example}<5->
        Let $B = \set{1,2,3}$ and $A = \set{a,b}$, then the collection $A^B$ is defined by the set of ordered tuples \EQ{\set{(aaa), (aab), (aba), (abb), (baa), (bab), (bba), (bbb)}}
    \end{example}    
\end{frame}

\begin{frame}{Inverse Map}
    \begin{block}{Definition}<1->
        For a function $f \in A^B$, we define \textbf{set inverse map} $f^{-1}(C) = \set{b \in B: f(b) \in C}$ for all subsets $C \subseteq A$. 
    \end{block}
    \begin{rem}<2-> 
        This is a slight abuse of notation since $f^{-1}: \cP(A) \to \cP(B)$ is a map from sets to sets. 
    \end{rem}
    \begin{example}<3->
        Let $B = \set{1,2,3}$ and $A = \set{a,b}$ and $f$ be denoted by the ordered tuple $(aba)$, then $f^{-1}(\set{a}) = \set{1,3}$ and $f^{-1}(\set{b}) = \set{2}$.
    \end{example}
\end{frame}

\begin{frame}{Injective, Surjective, Bijective} 
    A function $f \in A^B$ is
    \begin{itemize}
        \item<1->[\textbf{injective:}] if for any distinct $b \neq c \in B$, we have $f(b) \neq f(c)$,
        \item<2->[\textbf{surjective:}] if $f(B) = A$, and 
        \item<3->[\textbf{bijective:}] if it is injective and surjective.  
    \end{itemize}
\end{frame}

\begin{frame}{Examples}
    \begin{block}{Injective}
        Let $B = \set{1,2,3}$ and $A = \set{a,b,c,d}$. Then $(abc)$ is an injective function.   
    \end{block}
    \begin{block}{Surjective}<2->
        Let $B = \set{1,2,3,4}$ and $A = \set{a,b,c}$. Then $(abca)$ is a surjective function.
    \end{block}
    \begin{block}{Bijective}<3->
            Let $B = \set{1,2,3}$ and $A = \set{a,b,c}$. Then $(abc)$ is a bijective function.    
    \end{block}   
\end{frame}

\begin{frame}{Cardinality}
    \begin{block}{Definition}
        \begin{itemize}
            \item We denote the cardinality of a set $A$ by $\abs{A}$.
            \item<2->Cardinality describes a relationship between sets which compares their relative size.
            \item<3->If there is a bijection between two sets, they have the same cardinality.\\
            \item<4->Any set which is bijective to the set $[N]$ has cardinality $N$.
        \end{itemize}
    \end{block}

    \begin{Example}<5->
        The cardinality of $A = \set{a,b,c}$ is $\abs{A} = 3$, since there is a bijection between $B = \set{1,2,3}$ and $A = \set{a,b,c}$.  
    \end{Example}
\end{frame}

\begin{frame}{Countable}
    \begin{block}{Definition}
        \begin{itemize}
            \item<1-> Any set which is bijective to a subset of natural numbers $\N$ is called a \textbf{countable} set.
            \item<2->Any set which has a finite cardinality is called a \textbf{countably finite} set.
            \item<3->Any set which is bijective to the set of natural numbers $\N$ is called a \textbf{countably infinite} set.
        \end{itemize}
    \end{block}
\end{frame}

\begin{frame}{Exercise}
    \begin{block}{Show that}
        \begin{enumerate}
            \item $\abs{A^B} = \abs{A}^{\abs{B}}$. 
            \item $A^{[N]}$ is set of all $A$-valued $N$-length sequences. 
            \item $A^\N$ is a set of all $A$-valued countably infinite sequences indexed by the set of natural numbers $\N$.
            \item The sets $\N, \Z_+, \Z, \Q$ have the same cardinality. 
        \end{enumerate}
    \end{block}
\end{frame}

\begin{frame}{Sample Space} 
    \begin{block}{Definition}<1->
        The set of all possible outcomes of a random experiment is called \textbf{sample space} and denoted by $\Omega$. 
    \end{block}
    \begin{Example}[Single coin toss]<2->
        Consider a single coin toss, where the outputs can be heads or tails, denoted by $H$ and $T$, respectively.
        The set of all possible outcomes is $\Omega = \set{H, T}$.
    \end{Example}
    \begin{Example}[Finite coin tosses]<3->
    Consider $N$ finite tosses of a single coin. In this case, the sample space is $\Omega = \set{H, T}^{[N]}$, and a single outcome is a sequence $\omega = (\omega_1, \dots, \omega_N)$ where $\omega_i \in \set{H,T}$ for all $i \in [N]$.  
    \end{Example}
\end{frame}
\begin{frame}{Examples}
    \begin{Example}[Countably infinite coin tosses]<1->
        Consider an infinite sequence of coin tosses. 
        The set of all possible outcomes is $\Omega = \set{H, T}^\N$. 
        A single outcome is a sequence $\omega = (\omega_i \in \set{H,T}: i \in \N) \in \Omega$. 
    \end{Example}

    \begin{Example}[Point on non-negative real line]<2->
        The set of all possible outcomes for a single point on non-negative real line is $\Omega = \R_+$. 
    \end{Example}
    \begin{Example}[Countable points on non-negative real line]<3->
        The set of all possible outcomes is $\Omega = \R_+^\N$, where a single outcome is a sequence $\omega = (\omega_i \in \R_+: i \in \N) \in \Omega$. 
    \end{Example}
\end{frame}


\begin{frame}{Event space}
    \begin{block}{Definition}<1->
    A collection of subsets of sample space $\Omega$ is called the \textbf{event space} if it is a $\sigma$-algebra over subsets of $\Omega$, and denoted by $\sF$. 
    \end{block}
    \begin{block}{Properties}<2->
        In other words, the collection $\sF$ satisfies the following properties. 
        \begin{enumerate}
            \item<3-> Event space includes the certain event $\Omega$. That is, $\Omega \in \sF$. 
            \item<4-> Event space is closed under complements. That is, if $A \in \sF$, then $A^c \in \sF$. 
            \item<5-> Event space is closed under countable unions. That is, if $A_i \in \sF$ for all $i \in \N$, then $\cup_{i \in \N}A_i \in \sF$.
        \end{enumerate}
    \end{block}
    \begin{block}{Note}<6->
        The elements of the event space $\sF$ are called \textbf{events}. 
    \end{block}
\end{frame}

\begin{frame}{Examples}
    \begin{Example}[Coarsest event space]<1->
        For any sample space $\Omega$, the trivial event space is $\sF = \set{\emptyset, \Omega}$. 
    \end{Example}
    \begin{Example}[Single coin toss]<2-> 
        Recall that the sample space for a single coin toss is $\Omega = \set{H, T}$. We define an event space $\sF \triangleq \set{\emptyset, \set{H}, \set{T}, \set{H,T}}$. Can you verify that $\sF$ is a $\sigma$-algebra?
    \end{Example}
    \begin{Example}[Finest event space for finite coin tosses]<3-> 
        Recall that the sample space for $N$ coin tosses is $\Omega = \set{H, T}^N$, 
        %Consider $N$ tosses of a single coin, where the possible output of each coin toss belongs to the set $\set{H,T}$ as before.  
        %In this case, the sample space is 
        %\EQ{
        %\Omega = \set{H, T}^{[N]} = \set{(\omega_1, \dots, \omega_N): \omega_i \in \set{H,T} \text{ for all } i \in [N]}.
        %} 
        An event space for this sample space is $\sF \triangleq \cP(\Omega) = \set{A: A \subseteq \Omega}$. Can you verify that $\sF$ is a $\sigma$-algebra?
    \end{Example}
\end{frame}

\begin{frame}
    \begin{block}{Remark}
        %From the definition of an event space, the following statements hold true.
        %\begin{enumerate}
        %\item 
        We observe that $\sF \subseteq \cP(\Omega)$. 
        However, not all subsets of the sample space $\Omega$ necessary belong to the event space. 
        That is, $\sF \subset \cP(\Omega)$. 
        To show this, it suffices to construct a set $A \subset \Omega$ that is not in $\sF$. 
    \end{block}
\end{frame}

\begin{frame}{Properties of event space}
    \begin{block}{Property}<1->
        Event space contains the impossible event, is closed under finite unions, and closed under countable intersections.
    \end{block}
    \begin{Proof}<2->
        \begin{itemize}
            \item<2-> From the inclusion of certain event and the closure complements of $\sigma$-algebras, it follows that the impossible event $\emptyset \in \sF$. 
            
            \item<2-> That is, since $\Omega \in \sF$ and $\emptyset = \Omega^c$, we have $\emptyset \in \sF$.
            
            \item<3-> We need to show that if $A,B \in \sF$, then the union $A\cup B \in \sF$. This follows from the fact that we can $A_1 = A, A_2 = B$ and $A_i = \emptyset$ for all $i \ge 2$, and apply closure under countable unions. 
        \end{itemize}
    \end{Proof}
\end{frame}

\begin{frame}{Proof continued}
    \begin{Proof}
        \begin{itemize}
            \item<1-> We need to show that for any sequence of events such that $A_i \in \sF$ for all $i \in \N$, we have $\cap_{i \in \N}A_i \in \sF$. 
        
            \item<2-> To show this, we first notice that $A_i^c \in \sF$ for all $i \in \N$ from the closure under complements. 
        
            \item<3-> Further, we have $\cup_{i \in \N}A_i^c \in \sF$ from the closure under countable unions, and then the result follows from taking the complement and the closure under complements. 
        \end{itemize}
    \end{Proof}
\end{frame}
    
%\begin{rem}
%\item 
%\end{rem}
%\begin{rem}
%\item 
%Event space is closed under finite unions. 
%In particular, if $A, B \in \sF$, then $A\cup B \in \sF$. 

%\end{enumerate}
%\end{rem}
%\begin{rem}
%\item 
%Event space is closed under countable intersections. 
%That is, $A_i \in \sF$ for all $i \in \N$, then $\cap_{i \in \N}A_i \in \sF$. 
%\end{enumerate}
%\end{rem}


\begin{frame}{Event space generated by a family of sets}
    \begin{Definition}[Event space generated by a family of sets]<1->
        Consider a sample space $\Omega$ and a family $F \subset \cP(\Omega)$ of subsets of $\Omega$.   
        Then the event space generated by $F$ is the smallest event space containing each element of $F$, and is denoted by $\sigma(F)$. 
    \end{Definition}
        
    \begin{Example}[Event space generated by a single event]<2->
        For any sample space $\Omega$ and a subset $A \subseteq \Omega$, the smallest event space generated by this event $A$ is $\sF = \set{\emptyset, A, A^c, \Omega}$. 
    \end{Example}
\end{frame}

\begin{frame}{Examples continued}
    \begin{Example}[Countably infinite coin tosses] 
        Recall that the sample space for a countably infinite number of coin tosses is $\Omega = \set{H, T}^\N$.  
        %A single coin is tossed countably infinite times. 
        %A single outcome of this repeated coin toss can be denoted by a sequence $\omega = (\omega_1, \omega_2, \dots)$, 
        %where $\omega_i \in \set{H, T}$ is the output of $i$th coin toss. 
        %Then, we can write the sample space as 
        %\EQ{
        %\Omega = \set{H, T}^{[N]} = \set{(\omega_1, \omega_2, \dots ): \omega_i \in \set{H,T} \text{ for all } i \in \N}.
        %} 
        We will construct an event space $\sF$ on this outcome space, which would not be the power set of the sample space. 
        By definition, we must have the certain and the impossible event in an event space.
        We consider the event space $\sF$ generated by the events 
        \EQN{\label{eqn:Axioms:EventSpace:OneHeadEvent}
        A_n \triangleq \set{\omega \in \Omega: \omega_i = H \text{ for some }i \in [n]},\text{ for each }n \in \N.} 
        That is $A_n$ is the event of getting at least one head in first $n$ tosses. 
        We see that $(A_n \in \sF: n \in \N)$ is a sequence of increasing events.  
        From the closure under countable union, we have $\cup_{n \in \N}A_n = \Omega\setminus(T, T, \dots) \in \sF$. 
% and hence $(B_n \triangleq A_{n+1}\setminus A_n \in \sF: n \in \N)$ is a sequence of disjoint events, 
%where the event $B_n$ corresponds to seeing first head after $n-1$ consecutive tails. 
    \end{Example}
\end{frame}


\begin{frame}{Exercise}
    \begin{block}{Try yourself:}
        Consider a countably infinite sequence of coin tosses, with the sample space $\Omega = \set{H,T}^\N$ and the the event space $\sF$ generated by the events $(A_n: n \in \N)$ defined in Eq.~\eqref{eqn:Axioms:EventSpace:OneHeadEvent}. 
        Let $B_n$ be the event of observing first head in $n$th toss, 
        then show that $B_n \in \sF$ for all $ n \in \N$.  
    \end{block}
\end{frame}

%\begin{proof}
%We see that the event $B_1 = A_1$ and the event $B_{n+1}$ can be written as 
%\EQ{
%B_{n+1} = \set{\omega \in \Omega: \omega_i = T \text{ for }i \le n, \omega_{n+1} = H} = A_{n+1}\setminus A_{n}. 
%}
%Therefore, $B_n \in \sF$ for all $n \in \N$. 
%\end{proof}

\begin{frame}{Borel event spaces}
    \begin{Definition}<1->
        For sample space $\R$, a \emph{Borel event space} is generated by the subsets $B_x \triangleq (-\infty, x] \subseteq \R$ for each $x \in \R$, and denoted by $\sB(\R) = \sigma(\set{B_x: x \in \R})$. 
    \end{Definition}

    \begin{block}{Show the following}<2->
        The events $\set{x}, (x, y), [x, y]$ belong to the Borel event space $\sB(\R)$ for all $x, y \in \R$.  
    \end{block}
\end{frame}

\begin{frame}{Solution}
    \begin{itemize} 
        \item<1-> Let $x,y\in\R$ and $x < y$ , then we will show that the events $\set{x}, (x, y), [x, y]$ belong to Borel event space $\sB(\R)$.  
        \item<2-> We observe that $(-\infty, x]^c = (x,\infty) \in \sB(\R)$. 
        \item<3-> Similarly, $(y, \infty) \in \sB(\R)$. 
        \item<4-> It follows that $(x,\infty)\cap(-\infty, y] = (x, y] \in \sB(\R)$. 
        \item<5-> Consider a decreasing sequence $u$ such that $u_n \downarrow x$ and an increasing sequence $\ell$ such that $\ell_n\uparrow x$. 
        \item<6-> We observe that $(-\infty, \ell_n]\in \sB(\R)$ and $\cup_{n\in\N}(-\infty, \ell_n] = (-\infty,x) \in \sB(\R)$.  
        \item<7-> Therefore, it follows that $[x,y] = (-\infty, x)^c\cap(-\infty, y] \in \sB(\R)$. 
        \item<8-> We can similarly show that $(-\infty, y)\in \sB(\R)$, and it follows that $(x,y) = (-\infty, x]^c\cap(-\infty, y) \in \sB(\R)$. 
        \item<9-> We observe that $x \in [\ell_n,u_n] \in \sB(\R)$ for all $n \in \N$, and $\set{x} = \cap_{n \in \N}[\ell_n, u_n] \in \sB(\R)$. 
    \end{itemize}
\end{frame}



\end{document}